\documentclass[12pt,a4paper]{article}

\usepackage[utf8]{inputenc}
\usepackage[T1]{fontenc}
\usepackage[spanish]{babel}
\usepackage[top=2.5cm, bottom=2.5cm, left=2.5cm, right=2.5cm]{geometry}
\usepackage{booktabs}
\usepackage{longtable}
\usepackage{array}
\usepackage{listings}
\usepackage{fancyvrb}
\usepackage{xcolor}
\usepackage{parskip}

\definecolor{codigofondo}{RGB}{248,248,248}

\lstset{
  backgroundcolor=\color{codigofondo},
  basicstyle=\ttfamily\small,
  frame=single,
  breaklines=true,
  showstringspaces=false,
  numbers=left,
  numberstyle=\tiny\color{gray},
  numbersep=8pt,
}

\begin{document}

% ── Portada ──────────────────────────────────────────────────────────────────
\begin{center}
  {\LARGE\bfseries Proyecto Final --- Diseño de Mini-Lenguaje de Programación}\\[1em]
  {\large\bfseries Champions / CR7}\\[1.5em]
  \begin{tabular}{l}
    Juan Pablo Olave Muñoz \\
    Juan Sebastian López Guzmán \\
    Cristian Camilo Salazar Arenas \\
  \end{tabular}\\[0.5em]
  Procesadores del Lenguaje --- Semestre 5 --- 2025
\end{center}

\bigskip
\hrule
\bigskip

% ════════════════════════════════════════════════════════════════════════════
\section*{1. Análisis del lenguaje base asignado: OCaml}
% ════════════════════════════════════════════════════════════════════════════

\subsection*{1.1 Componentes léxicos}

\textbf{Palabras reservadas:}
\texttt{let}, \texttt{in}, \texttt{fun}, \texttt{if}, \texttt{then}, \texttt{else},
\texttt{while}, \texttt{do}, \texttt{done}, \texttt{for}, \texttt{to}, \texttt{downto},
\texttt{begin}, \texttt{end}, \texttt{match}, \texttt{with}, \texttt{rec},
\texttt{true}, \texttt{false}, \texttt{try}, \texttt{raise}, \texttt{exception},
\texttt{type}, \texttt{module}, \texttt{open}, \texttt{mutable}, \texttt{ref}

\textbf{Operadores:}
\begin{itemize}
  \item Aritméticos: \texttt{+}, \texttt{-}, \texttt{*}, \texttt{/}, \texttt{mod}
  \item Relacionales: \texttt{=}, \texttt{<>}, \texttt{<}, \texttt{>}, \texttt{<=}, \texttt{>=}
  \item Lógicos: \texttt{\&\&}, \texttt{||}, \texttt{not}
  \item Asignación de referencia: \texttt{:=}, \texttt{<-}
  \item Concatenación de cadenas: \texttt{\^{}}
\end{itemize}

\textbf{Delimitadores:}
\texttt{(}, \texttt{)}, \texttt{[}, \texttt{]}, \texttt{\{}, \texttt{\}},
\texttt{;;}, \texttt{;}, \texttt{->}, \texttt{|}, \texttt{::}

\textbf{Tipos de datos primitivos:}
\texttt{int}, \texttt{float}, \texttt{string}, \texttt{bool}, \texttt{char}, \texttt{unit}

\textbf{Estructuras de control:}
\begin{itemize}
  \item Condicional: \texttt{if} \ldots \texttt{then} \ldots \texttt{else}
  \item Ciclo while: \texttt{while} \ldots \texttt{do} \ldots \texttt{done}
  \item Ciclo for: \texttt{for} \textit{var} \texttt{=} \textit{inicio} \texttt{to} \textit{fin} \texttt{do} \ldots \texttt{done}
  \item Pattern matching: \texttt{match} \ldots \texttt{with} \texttt{|} \ldots \texttt{->} \ldots
\end{itemize}

\subsection*{1.2 Tabla léxica del lenguaje base}

\begin{center}
\renewcommand{\arraystretch}{1.25}
\begin{longtable}{>{\ttfamily}p{2.8cm} p{3.5cm} >{\ttfamily}p{4cm} p{3.5cm}}
\toprule
\normalfont\textbf{Token} & \textbf{Categoría} & \normalfont\textbf{Patrón (regex)} & \textbf{Equivalente} \\
\midrule
\endfirsthead
\toprule
\normalfont\textbf{Token} & \textbf{Categoría} & \normalfont\textbf{Patrón (regex)} & \textbf{Equivalente} \\
\midrule
\endhead
\bottomrule
\endfoot

let        & Palabra reservada    & let                          & Declaración de variable/función \\
in         & Palabra reservada    & in                           & Alcance de un \texttt{let} \\
if         & Palabra reservada    & if                           & Condicional \\
then       & Palabra reservada    & then                         & Rama verdadera \\
else       & Palabra reservada    & else                         & Rama falsa \\
while      & Palabra reservada    & while                        & Ciclo while \\
do         & Palabra reservada    & do                           & Inicio del cuerpo del ciclo \\
done       & Palabra reservada    & done                         & Fin del cuerpo del ciclo \\
for        & Palabra reservada    & for                          & Ciclo for \\
match      & Palabra reservada    & match                        & Pattern matching \\
true/false & Literal booleano     & true|false                   & Valor booleano \\
rec        & Palabra reservada    & rec                          & Función recursiva \\
=          & Op. igualdad         & =                            & Igualdad / asignación \\
<>         & Op. relacional       & <>                           & Desigualdad \\
< >        & Op. relacional       & [<>]                         & Menor / mayor que \\
<= >=      & Op. relacional       & <=|>=                        & Menor o igual / mayor o igual \\
+ - * /    & Op. aritméticos      & [+\textbackslash-*/]         & Suma, resta, multiplicación, división \\
mod        & Op. aritmético       & mod                          & Módulo (resto) \\
\&\&       & Op. lógico           & \&\&                         & AND lógico \\
||         & Op. lógico           & ||                           & OR lógico \\
not        & Op. lógico           & not                          & Negación lógica \\
\textasciicircum & Op. concatenación & \textasciicircum           & Concatenar cadenas \\
::         & Op. de lista         & ::                           & Constructor cons \\
( )        & Delimitador          & [()]                         & Agrupación \\
[ ]        & Delimitador          & [\textbackslash[\textbackslash]] & Lista \\
;;         & Delimitador          & ;;                           & Fin de expresión top-level \\
->         & Delimitador          & ->                           & Flecha (match / tipo función) \\
INTEGER    & Literal entero       & [0-9]+                       & Número entero \\
FLOAT      & Literal flotante     & [0-9]+\textbackslash.[0-9]* & Número flotante \\
STRING     & Literal de cadena    & "[{\^{}}"]*"                 & Cadena de texto \\
CHAR       & Literal de carácter  & '[{\^{}}'\\]'                & Carácter individual \\
IDENT      & Identificador        & [a-z\_][a-zA-Z0-9\_']*      & Nombre de variable o función \\

\end{longtable}
\end{center}

% ════════════════════════════════════════════════════════════════════════════
\section*{2. Diseño del nuevo lenguaje}
% ════════════════════════════════════════════════════════════════════════════

\subsection*{2.1 Nombre, temática e inspiración}

\textbf{Nombre:} Champions (archivos \texttt{.cr7})

\textbf{Temática:} Fútbol moderno, inspirado en Cristiano Ronaldo (CR7), el Champions League y la jerga futbolística.

\textbf{Justificación:} El equipo eligió el fútbol porque permite mapear los conceptos de programación de forma intuitiva: declarar una variable es ``pasar el balón'', un condicional es una ``jugada de ataque'', un ciclo es una ``repetición'', imprimir es ``gritar gol''. Cada palabra clave tiene una razón futbolística clara, lo que le da identidad y coherencia al lenguaje.

\subsection*{2.2 Tabla de equivalencias léxicas}

\begin{center}
\renewcommand{\arraystretch}{1.3}
\begin{longtable}{>{\ttfamily}p{4cm} p{3cm} p{6.5cm}}
\toprule
\normalfont\textbf{Palabra en Champions} & \textbf{Equivalente OCaml} & \textbf{Inspiración / temática} \\
\midrule
\endfirsthead
\toprule
\normalfont\textbf{Palabra en Champions} & \textbf{Equivalente OCaml} & \textbf{Inspiración / temática} \\
\midrule
\endhead
\bottomrule
\endfoot

pasar\_balon      & \texttt{let}        & Pasar el balón a un compañero — entregar un valor a una variable \\
a                 & \texttt{=} (asig.)  & Preposición de dirección: ``pasar el balón \textbf{a} x'' \\
gritar\_gol       & \texttt{print}      & Cuando hay gol se grita — mostrar un valor en pantalla \\
ataca\_cr7        & \texttt{if}         & CR7 ataca cuando la condición es verdadera \\
defiende\_cr7     & \texttt{else if}    & Alternativa si la primera condición falló \\
amaga\_cr7        & \texttt{else}       & Amague final cuando ninguna condición fue verdadera \\
ataca\_el\_atleti & \texttt{while}      & El Atlético defiende en ciclos interminables — bucle con condición \\
ataca\_el\_arsenal & \texttt{for}       & El Arsenal ataca con estructura definida: inicio, condición y actualización \\
ofensiva          & \texttt{++}         & Subir la línea ofensiva — incrementar en 1 \\
defensiva         & \texttt{--}         & Bajar la línea defensiva — decrementar en 1 \\
gol               & \texttt{true}       & El gol es el resultado positivo — condición verdadera \\
autogol           & \texttt{false}      & El autogol es lo contrario — condición falsa \\
cabeceo           & \texttt{+}          & Cabecear suma altura al salto — suma aritmética \\
rabona            & \texttt{-}          & La rabona cruza las piernas — resta \\
chilena           & \texttt{*}          & La chilena multiplica el espectáculo — multiplicación \\
penalti           & \texttt{/}          & El penalti divide el juego en dos momentos — división \\
volea             & \texttt{mod}        & La volea devuelve el balón al borde — módulo / resto \\
mbappe            & \texttt{\^{}}       & Mbappé conecta jugadas con velocidad — concatenar cadenas \\
tikitaka          & \texttt{+=}         & El tiki-taka acumula pases — sumar y reasignar \\
mudryk            & \texttt{-=}         & Mudryk pierde el balón — restar y reasignar \\
transfermarkt     & \texttt{==}         & Transfermarkt cotiza si dos jugadores valen igual — igualdad \\
vendehumo         & \texttt{!=}         & El vendehumo promete algo que no es — desigualdad \\
barredora         & \texttt{>}          & La entrada en barredora supera al rival — mayor que \\
amarilla          & \texttt{>=}         & Tarjeta amarilla: límite alcanzado o superado — mayor o igual \\
roja              & \texttt{<}          & Tarjeta roja: por debajo del límite — menor que \\
zancadilla        & \texttt{<=}         & La zancadilla lo tumba antes del límite — menor o igual \\
pase              & \texttt{\&\&}       & El pase requiere que ambos jugadores estén listos — AND lógico \\
sombrero          & \texttt{||}         & El sombrero engaña: con una de dos opciones basta — OR lógico \\
\texttt{!}        & \texttt{;;}         & El silbato del árbitro termina la jugada — fin de sentencia \\
\texttt{\{ \}}    & \texttt{begin/end}  & Las porterías delimitan el campo — inicio y fin de bloque \\

\end{longtable}
\end{center}

\subsection*{2.3 Gramática básica (BNF simplificado)}

\begin{Verbatim}[frame=single, fontsize=\small]
<programa>       ::= <sentencia>+

<sentencia>      ::= <declaracion>
                   | <asignacion>
                   | <asig_compuesta>
                   | <incremento>
                   | <impresion>
                   | <condicional>
                   | <ciclo_while>
                   | <ciclo_for>
                   | <bloque>

<declaracion>    ::= 'pasar_balon' IDENT 'a' <expresion> '!'

<asignacion>     ::= IDENT 'a' <expresion> '!'

<asig_compuesta> ::= IDENT 'tikitaka' <expresion> '!'
                   | IDENT 'mudryk'   <expresion> '!'

<incremento>     ::= IDENT 'ofensiva' '!'
                   | IDENT 'defensiva' '!'

<impresion>      ::= 'gritar_gol' '(' <expresion> ')' '!'

<condicional>    ::= 'ataca_cr7' '(' <expresion> ')' <bloque>
                     ['defiende_cr7' '(' <expresion> ')' <bloque>]*
                     ['amaga_cr7' <bloque>]

<ciclo_while>    ::= 'ataca_el_atleti' '(' <expresion> ')' <bloque>

<ciclo_for>      ::= 'ataca_el_arsenal'
                     '(' 'pasar_balon' IDENT 'a' <expresion> '!'
                         <expresion> '!'
                         IDENT ('ofensiva' | 'defensiva') ')'
                     <bloque>

<bloque>         ::= '{' <sentencia>* '}'

<expresion>      ::= <expr_or>
<expr_or>        ::= <expr_and> ('sombrero' <expr_and>)*
<expr_and>       ::= <expr_cmp> ('pase' <expr_cmp>)*
<expr_cmp>       ::= <expr_add> [<op_cmp> <expr_add>]
<expr_add>       ::= <expr_mul> (('cabeceo'|'rabona'|'mbappe') <expr_mul>)*
<expr_mul>       ::= <factor> (('chilena'|'penalti'|'volea') <factor>)*
<factor>         ::= INT_LIT | STR_LIT | BOOL_LIT | IDENT | '(' <expresion> ')'

<op_cmp>         ::= 'transfermarkt' | 'vendehumo' | 'roja' | 'barredora'
                   | 'zancadilla' | 'amarilla'

BOOL_LIT         ::= 'gol' | 'autogol'
INT_LIT          ::= [0-9]+
STR_LIT          ::= '"' [^"]* '"'
IDENT            ::= [a-zA-Z_][a-zA-Z0-9_]*
\end{Verbatim}

% ════════════════════════════════════════════════════════════════════════════
\section*{3. Programa de prueba}
% ════════════════════════════════════════════════════════════════════════════

El siguiente archivo \texttt{basicos.cr7} ejercita al menos una declaración, un condicional y un ciclo:

\begin{lstlisting}[caption={examples/basicos.cr7}]
# --- 1. Declaracion de variables ---
pasar_balon mensaje a "Bienvenido al Champions"!
pasar_balon balon a 100!
pasar_balon contador a 0!

# --- 2. Impresion basica ---
gritar_gol(mensaje)!

# --- 3. Condicional if / else-if / else ---
ataca_cr7(balon zancadilla 7){
    gritar_gol("Mete gol CR7")!
}defiende_cr7(balon amarilla 10){
    gritar_gol("Defiende el arco")!
}amaga_cr7{
    gritar_gol("Gol de Mudryk")!
}

# --- 4. Ciclo while ---
pasar_balon i a 0!
ataca_el_atleti(i roja 5){
    gritar_gol("Repeticion " mbappe i)!
    i tikitaka 1!
}

# --- 5. Ciclo for ---
ataca_el_arsenal(pasar_balon j a 0! j roja 3! j ofensiva){
    gritar_gol("Indice " mbappe j)!
}

# --- 6. Operaciones aritmeticas ---
pasar_balon num1 a 20!
pasar_balon num2 a 5!
pasar_balon suma     a num1 cabeceo num2!
pasar_balon resta    a num1 rabona  num2!
pasar_balon producto a num1 chilena num2!
pasar_balon cociente a num1 penalti num2!
pasar_balon modulo   a num1 volea   num2!
gritar_gol("Suma: "     mbappe suma)!
gritar_gol("Resta: "    mbappe resta)!
gritar_gol("Producto: " mbappe producto)!
gritar_gol("Cociente: " mbappe cociente)!
gritar_gol("Modulo: "   mbappe modulo)!
\end{lstlisting}

\textbf{Salida esperada:}

\begin{Verbatim}[frame=single, fontsize=\small]
Bienvenido al Champions
Gol de Mudryk
Repeticion 0
Repeticion 1
Repeticion 2
Repeticion 3
Repeticion 4
Indice 0
Indice 1
Indice 2
Suma: 25
Resta: 15
Producto: 100
Cociente: 4
Modulo: 0
\end{Verbatim}

\end{document}
